\chapter{Conclusioni e sviluppi futuri}
\label{cap:conclusioni}

\section{Risposta alla domanda di partenza}

La domanda con cui si apre questa tesi era: \emph{il radar mmWave può sostituire o
affiancare il sensore PIR nel nodo DIPME-DEVICE del progetto UPRISE/SAFE?}

La risposta, sulla base dei dati raccolti, è \textbf{affiancare, non sostituire}, e
si articola in tre affermazioni sostenute da misure.

\paragraph{1. Sul caso d'uso che interessa il progetto, il PIR non funziona --- e
non per difetto ma per costruzione.} Con un soggetto immobile a 2,30~m, su cinque
trial da cinque minuti, il PIR ha mancato la presenza nel
$99{,}78 \pm 0{,}49\,\%$ dei campioni, con zero rilevamenti in quattro trial su
cinque; il radar l'ha rilevata nel $100\,\%$ dei 7554 campioni
(\cref{sec:persona-ferma}). Il risultato più significativo è però un altro: il PIR
fallisce anche con il soggetto \emph{in movimento sul posto}, con
$80{,}5\,\%$ di falsi negativi a 1~m e $100\,\%$ oltre i 2~m
(\cref{sec:pir-movimento}). Poiché una persona intrappolata si muove sul posto e non
attraversa la stanza, l'ipotesi implicita ``il PIR non vede chi è fermo ma vede chi
si muove'' è, nello scenario di interesse, falsa. Questa è la conclusione più
rilevante per il progetto.

\paragraph{2. Il radar non è solo più affidabile: fornisce grandezze diverse.} La
distanza è misurata con linearità eccellente ($R^2 = 0{,}99965$ su 25 trial fra 1 e
5~m) e con un errore riducibile sotto i 5~cm applicando un singolo coefficiente
correttivo (\cref{sec:distanza}). Le venti serie di energia dell'engineering mode
consentono di estrarre il ritmo respiratorio di un soggetto immobile, con
concordanza fra canali indipendenti (\cref{sec:respiro}), e sono la base
dell'indice di vitalità (\cref{cap:vitalita}). Nessuna di queste grandezze è
rappresentabile con l'uscita a un bit di un PIR. E il vantaggio non si paga in
affidabilità: su 7~ore di stanza vuota entrambi i sensori hanno prodotto zero falsi
positivi, con limite superiore al 95\,\% di $\leq 0{,}43$ eventi/h
(\cref{sec:falsi-positivi}).

\paragraph{3. Il consumo esclude la sostituzione secca e indica l'architettura.} Il
radar assorbe 80~mA contro i meno di 50~\textmu A del PIR: tre ordini di grandezza
(\cref{cap:consumi}). Un nodo integrato in un arredo deve restare in attesa per anni
a batteria, e con il radar sempre acceso l'autonomia si misura in ore. La soluzione
non è scegliere fra i due sensori ma metterli in sequenza temporale, secondo
l'architettura ibrida di \cref{sec:architettura-ibrida}: il PIR come guardiano a
microampere in ``tempo di pace'', il radar mmWave come strumento di misura durante la
finestra di emergenza, con \emph{duty-cycling} per estendere l'autonomia oltre le 72
ore. Il PIR non è un concorrente del radar: è ciò che decide quando accenderlo.

\medskip
\noindent
La raccomandazione operativa per il progetto è quindi di \textbf{aggiungere} un
modulo mmWave al DIPME-DEVICE, mantenendo il PIR nel ruolo di sorgente di risveglio,
e di attivare il radar sul trigger di modalità terremoto.

\section{Contributi consolidati}

Riassumendo quanto prodotto:

\begin{itemize}[leftmargin=1.4em,itemsep=3pt]
  \item una piattaforma di misura riproducibile su ESP32 che campiona radar e PIR in
        modo sincrono a 5,00~Hz esatti con accesso ai dati per-gate, con un formato
        CSV unico condiviso da firmware, acquisizione, analisi e interfaccia web;
  \item la caratterizzazione quantitativa, con cinque ripetizioni per scenario, del
        divario PIR/mmWave sulla persona ferma e sulla persona in movimento sul
        posto;
  \item la calibrazione della misura di distanza del LD2410B e l'individuazione di un
        errore di scala sistematico correggibile con un coefficiente;
  \item la dimostrazione di misurabilità del respiro con hardware consumer, con
        l'individuazione documentata di quali canali siano effettivamente
        utilizzabili e perché gli altri non lo siano;
  \item la specifica di un indice di vitalità implementabile su microcontrollore, con
        un percorso di taratura e validazione su trial separati;
  \item l'analisi dei consumi e l'argomentazione dell'architettura ibrida.
\end{itemize}

\subsection{Risultati metodologici collaterali}

Tre esiti non previsti meritano una menzione, perché sono il tipo di informazione
che raramente entra nella documentazione ma che ha valore per chi ripete il lavoro.

\begin{enumerate}[leftmargin=1.8em,itemsep=3pt]
  \item \textbf{Non fidarsi degli strumenti di configurazione da PC per i valori
        numerici.} Lo strumento ufficiale \texttt{LD2410 Tool} riporta sensibilità
        pari a 100 su tutti i gate; la lettura via UART restituisce la scala di
        fabbrica decrescente, coerente con il protocollo e con il comportamento
        osservato (\cref{sec:identita}). Per i parametri va usata la lettura diretta
        dal modulo.
  \item \textbf{Fare riferimento alla tabella di piedinatura del datasheet, non
        all'ordine apparente dei fili.} Un'inversione dei due fili della seriale ha
        prodotto tre settimane di lavoro nella convinzione che il modulo fosse
        guasto (\cref{sec:cablaggio}).
  \item \textbf{Verificare la qualità dei dati prima di analizzarli.} La saturazione
        dei canali di energia, i canali stazionari strutturalmente nulli e la coda di
        presenza (\cref{sec:insidie}) sono tutti fenomeni che, non riconosciuti,
        producono conclusioni sbagliate anziché errori evidenti: un canale saturo non
        segnala il proprio stato, restituisce semplicemente numeri privi di
        informazione. Da qui la scelta di eseguire un controllo di qualità
        automatico su ogni file prima di usarlo.
\end{enumerate}

\section{Lavoro rimanente}

Per completezza si elencano le attività pianificate e non ancora completate alla
data di questa stesura, con la sezione in cui sono collocate:

\begin{itemize}[leftmargin=1.4em,itemsep=2pt]
  \item latenze di rilevamento e di rilascio con la convenzione dell'evento a
        istante noto (\cref{sec:latenza});
  \item scenario ``persona sotto il banco'' a distanza inferiore al metro
        (\cref{sec:sotto-banco});
  \item penetrazione degli ostacoli: cartongesso, legno, vetro, plastica
        (\cref{sec:ostacoli});
  \item prova del respiro con ritmo imposto da metronomo e controllo negativo su
        stanza vuota (\cref{sec:respiro-metronomo});
  \item riconfigurazione del LD2420 in modalità binaria per il confronto quantitativo
        fra i due radar (\cref{sec:ld2420-risultati});
  \item implementazione della dashboard e test di equivalenza dei due CSV
        (\cref{sec:webui-piano});
  \item taratura e validazione dell'indice di vitalità
        (\cref{sec:vitalita-taratura}).
\end{itemize}

\section{Sviluppi futuri}

\paragraph{Verso il dimostratore reale.} La campagna è stata svolta in ambiente
domestico con un solo soggetto. Il passo successivo naturale è la ripetizione degli
scenari chiave nei dimostratori del progetto --- il Polo Lodovici di Informatica e
l'IIS Fermi Sacconi Ceci di Ascoli Piceno --- dove dimensioni dell'ambiente,
materiali, riflessioni multiple e presenza di più persone sono realistici. Le
domande aperte sono due: il tasso di falsi positivi in un'aula con arredi metallici,
e il comportamento con più persone contemporaneamente presenti, che è il limite
strutturale del LD2410B (\cref{sec:discussione}).

\paragraph{Il montaggio sulla lamiera forata.} È la questione di progetto più
concreta e va posta ai progettisti dell'arredo. Il piano del banco è rinforzato con
lamiera forata antisfondamento, e il metallo è il materiale peggiore per un radar a
24~GHz. Tre configurazioni sono da valutare sperimentalmente: sensore montato
\emph{sotto} il piano e quindi rivolto verso il volume di rifugio (ipotesi
preferibile); sensore montato \emph{sopra} il piano, che richiede di verificare se il
diametro e il passo dei fori consentono la trasmissione a 24~GHz --- una griglia
metallica con fori molto più piccoli della lunghezza d'onda di 12,5~mm si comporta
come uno schermo continuo, quindi la geometria della perforazione è determinante;
sensore integrato nel telaio laterale. Questa verifica è un'estensione di
\cref{sec:ostacoli} con il materiale reale.

\paragraph{Selettività fra banchi adiacenti.} Nel banco UPRISE i banchi sono
strutturalmente connessi fra loro, e un radar con portata di 6~m vede diversi banchi.
Il problema si affronta limitando il gate massimo via UART, funzione già
implementata e verificata nel corredo software di questa tesi. La misura da fare è
il tasso di attribuzione corretta con due volumi di rifugio adiacenti, che è
l'informazione necessaria a una mappa di calore utile ai soccorritori: sapere che
c'è qualcuno in un'aula è molto meno utile che sapere sotto quale banco.

\paragraph{Integrazione con il nodo DIPME e la rete LoRa.} L'indice di vitalità è
progettato per essere calcolato a bordo e occupa pochi byte: è esattamente il tipo di
informazione trasmissibile su LoRa, dove il vincolo è la dimensione del messaggio e
il duty cycle. Un messaggio che riporti presenza, distanza, classe di vitalità e
istante dell'ultima rilevazione sta in pochi byte ed è direttamente rappresentabile
sulla mappa di calore dei soccorritori. L'integrazione con il DIPME-DEVICE e con il
Digital Twin del progetto~\cite{callisto2023dt} è il proseguimento naturale del
lavoro.

\paragraph{Analisi in frequenza a bordo.} L'estensione dell'indice di vitalità con
una FFT calcolata sull'ESP32 su finestra scorrevole di 30~s renderebbe il termine
respiratorio più robusto agli artefatti rispetto alla media esponenziale attuale, a
costo di un buffer e di qualche millisecondo di calcolo per finestra: entrambi
disponibili sull'hardware in uso.

\paragraph{Confronto diretto con l'UWB del DIPME.} Il nodo esistente monta già un
sensore UWB. Un confronto sperimentale a tre --- PIR, mmWave, UWB --- sullo stesso
scenario e con la stessa metodologia sarebbe il completamento naturale di questo
lavoro, e trasformerebbe l'argomentazione di \cref{sec:uwb}, oggi basata su
datasheet e letteratura, in un risultato misurato.
